\documentclass[10pt,a4paper]{amsart}
\usepackage[margin=19mm]{geometry}
\usepackage{amsmath,amssymb,mathtools}
\usepackage[scr=rsfs,cal=euler]{mathalfa}
\usepackage{xfrac}
\usepackage[unicode,hidelinks,bookmarks=false]{hyperref}
\newcommand{\ie}{i.e.\ }
\newcommand{\cf}{cf.\ }
\newcommand{\resp}{respectively\ }
\newcommand{\Subsection}{\subsection}
\providecommand{\nobreakdash}{\nobreak\hbox{-}}
\setlength{\emergencystretch}{2em}
\multlinegap0pt
\usepackage{bm}
\usepackage{bbm}
\usepackage{dsfont}
\usepackage{xspace}
\usepackage{cancel}
\usepackage{mathtools}
\usepackage[shortlabels]{enumitem}
\usepackage{amsthm}
\makeatletter
\@ifundefined{theorem}{\newtheorem{theorem}{Theorem}}{}
\@ifundefined{lemma}{\newtheorem{lemma}[theorem]{Lemma}}{}
\makeatother
\def\ZZ{\mathbb{Z}}
\newcommand{\OK}{\mathcal{O}_K}
\newcommand{\cbr}[1]{ \left\{ {#1} \right\} }
\newcommand{\isom}{\cong}
\title[Sharpness of the Mockenhaupt-Mitsis-Bak-Seeger Fourier restr]{Sharpness of the Mockenhaupt-Mitsis-Bak-Seeger Fourier restriction theorem in all dimensions}
\author{Robert Fraser, Kyle Hambrook,, Donggeun Ryou}
\date{Source: arXiv:2505.19526v1 (CC BY 4.0)}
\begin{document}
\maketitle

\noindent\emph{Source and licence.} Sharpness of the Mockenhaupt-Mitsis-Bak-Seeger Fourier restriction theorem in all dimensions. Version arXiv:2505.19526v1 (CC BY 4.0), distributed under the Creative Commons Attribution 4.0 International licence, as recorded in the arXiv licence metadata for this exact version and shown on the abstract page https://arxiv.org/abs/2505.19526v1. Full rights notice: licenses/arxiv-2505.19526-CC-BY-4.0.txt.\par\smallskip

\noindent\textit{Editorial scope.} Counted unit: the Lemma whose source label is \texttt{h\_2\_est} in the article (source0.tex lines 1843--1850, source label \texttt{h\_2\_est}), delivered together with the article's own complete original proof of it (source0.tex lines 1851--1941, one \texttt{proof} environment of 91 lines). No mathematics is written, repaired or re-derived: statement and proof are the article's own text line for line. Required context retained uncounted: idealquotient (lines 584--587); separation lemma (lines 616--623); concrete crt lemma (lines 716--727); Phi-Fourier-Lemma (lines 788--798); QMsize (lines 839--841). Every definition, display and label of those lines is retained unchanged. Omitted from this package and not redistributed: 1--583, 588--615, 624--715, 728--787, 799--838, 842--1842, 1942--2532. Nothing inside a retained excerpt is omitted or altered: every retained body is delivered exactly as the article's lines are written, so no omission receipt of any kind accompanies this package. Citations: the retained excerpts carry no citation command at all, so no bibliography entry of the article is transcribed and no reference is redistributed. Reference adaptation: none was needed and none was made; every reference inside the delivered unit resolves inside it, so no unresolved reference or citation is produced. Printed slips kept literal: no printed word, symbol, hypothesis or number of the counted statement or of its proof was altered. Layout and class: the article's own class is \texttt{amsart} with packages including ulem, enumitem, amssymb, amsmath, amsthm, verbatim, IEEEtrantools, hyperref, xcolor; the wrapper is the lane's pinned \texttt{amsart} editorial wrapper at 10pt on A4, which restates the theorem-like environments the retained text needs and adds the article's own macro definitions that the retained text needs (\texttt{\textbackslash OK}, \texttt{\textbackslash ZZ}, \texttt{\textbackslash cbr}, \texttt{\textbackslash isom}), with extra wrapper packages (bm, bbm, dsfont, xspace, cancel, mathtools, enumitem). The delivered numbering is the unit's own. Assets: no retained span contains a figure environment, a graphics-inclusion command, a PDF-page inclusion, a table or a TikZ picture; no image file is copied into this package, the package declares no asset dependency and no image file is redistributed with it. Licence: version arXiv:2505.19526v1 of this article is distributed under the Creative Commons Attribution 4.0 International licence, as recorded in the arXiv licence metadata for this exact version and shown on the abstract page https://arxiv.org/abs/2505.19526v1. A full rights notice is delivered at licenses/arxiv-2505.19526-CC-BY-4.0.txt. Characters: every retained excerpt is pure ASCII, so the delivered engine, XeLaTeX with its default Latin Modern text font, reports no missing character.
\begin{lemma}\label{idealquotient}
If $I_2 \subseteq I_1$ are fractional ideals of $\OK$,  
then  $N(I_2) = N(I_1)|I_1/I_2|$. 
\end{lemma}

\begin{lemma}\label{separation lemma}
Let $I_1$ and $I_2$ be non-zero fractional ideals of $\OK$. Let $r_1 \in I_1^{-1}$ and $r_2 \in I_2^{-1}$ with $r_1 \neq r_2$. Then 
$$
C_B|r_1 - r_2| 
\geq (N(I_1 \cap I_2))^{-1/d}
\geq (N(I_1) N(I_2))^{-1/d}. 
$$
\end{lemma}

\begin{lemma}\label{concrete crt lemma}
Suppose $I_1$, $I_2$, and $D$ are fractional ideals of $\OK$. 
Suppose $a_1, a_2 \in D$. 
\begin{enumerate}[(a)]
\item If $a_1 - a_2 \notin D(I_1 + I_2)$, then
\[(a_1 + DI_1) \cap (a_2 + DI_2) = \emptyset.\]
\item If $a_1 - a_2 \in D(I_1 + I_2)$, then there exists $a \in D$ such that 
$$
(a_1 + DI_1) \cap (a_2 + DI_2) = a + (D I_1 \cap DI_2).
$$
\end{enumerate}
\end{lemma}

\begin{lemma}\label{Phi-Fourier-Lemma}
Let $I$ be a non-zero ideal of $\OK$, let $s \in \ZZ^d \isom \OK$, and let 
$\eta > 0$. Then 
\begin{equation}\label{Phi-Fourier-eqn-1}
\widehat{\Phi_{I,\eta}}(s) = \widehat{\phi}(\eta s) N(I)\mathbf{1}_{\delta^{-1}I} (T^{-1} s). 
\end{equation}
Moreover, if $0 < \eta < C_B^{-1}N(I)^{-1/d}$, then  
\begin{equation}\label{Phi-Fourier-eqn-2}
\widehat{\Phi_{I,\eta}^2}(s) = \eta^{-d} \widehat{\phi^2}(\eta s) N(I)\mathbf{1}_{\delta^{-1}I} (T^{-1} s).
\end{equation}
\end{lemma}

\begin{equation}\label{QMsize}
|Q(M)| \sim \frac{M^d}{\log(M)}.
\end{equation}

\begin{lemma}\label{h_2_est}
Suppose $\rho \neq 0$. 
For all 
$s \in \ZZ^d$, 
$$
|\widehat{h_2}(s)| \lesssim  M_k^{d\tau} (1+ M_k^{\rho -d(\tau-1)})(1 + \eta_k |s|)^{-10d}.
$$
\end{lemma}

\begin{proof}
Note that
\begin{align*}
\widehat{h_2}(s) 
= \sum_{I \in Q(M_k)} \sum_{t \in \mathbb{Z}^d} 
c_k \widehat{\Phi_{I, \eta_k}}(t) \widehat{\Phi_{J_k, \eta_k}}(s - t).
\end{align*}
By Lemma \ref{Phi-Fourier-Lemma}, 
\[\widehat{h_2}(s) = \sum_{I \in Q(M_k)} \sum_{t \in \mathcal{T}(I)} 
c_k 
N(I) 
 N(J_k) 
 \widehat \phi(\eta_k t)
\widehat \phi(\eta_k(s - t)),\]
where 
$$
\mathcal{T}(I)
=
\cbr{t \in \ZZ^d : T^{-1}t \in \delta^{-1}I \cap (T^{-1}s + \delta^{-1}J_k)}.
$$
Since at least one of $|s - t|$ and $|t|$ will be $\geq |s|/2$, 
the Schwartz bounds on $\widehat{\phi}$ 
imply 
$$
|\widehat \phi(\eta_k t)
\widehat \phi(\eta_k(s - t))| 
\lesssim (1+\eta_k|s|)^{-10d}(1+\eta_k|t|)^{-50d}.
$$
Therefore,
\begin{equation}\label{h2hatsum2}
|\widehat{h_2}(s)| \lesssim 
(1 + \eta_k |s|)^{-10d} \log (M_k) M_k^{\rho} 
\sum_{I \in Q(M_k)} 
\sum_{t \in \mathcal{T}(I)} 
(1+\eta_k|t|)^{-50d}. 
\end{equation}
Consider a fixed $I \in Q(M_k)$. 
By Lemma \ref{concrete crt lemma},  either the set 
$\delta^{-1} I \cap (T^{-1} s + \delta^{-1} J_k)$ is empty or 
there exists $a \in \delta^{-1}$ such that 
$$
\delta^{-1} I \cap (T^{-1} s + \delta^{-1} J_k)
=
a + \delta^{-1} I \cap \delta^{-1}J_k. 
$$ 
In the first case, $\mathcal{T}(I)$ is empty, so the sum over $t \in \mathcal{T}(I)$ in \eqref{h2hatsum2} is zero. 
In the second case, 
$\delta^{-1} I \cap (T^{-1} s + \delta^{-1} J_k)$ is a translate of 
$\delta^{-1} I \cap \delta^{-1}J_k$, 
so the separation between the elements of $\delta^{-1} I \cap (T^{-1} s + \delta^{-1} J_k)$
 is the same as the separation between the elements of $\delta^{-1} I \cap \delta^{-1}J_k$. 
By Lemma \ref{separation lemma} (with $I_1 = I_2 = ( \delta^{-1} I \cap \delta^{-1}J_k )^{-1}$), the elements of $\delta^{-1} I \cap \delta^{-1}J_k$ are separated by 
$
\gtrsim 
N(\delta^{-1} I \cap \delta^{-1}J_k)^{1/d}.
$ 
But, by Lemma \ref{idealquotient}, 
\begin{align}\label{I Jk norm calc}
N(\delta^{-1} I \cap \delta^{-1}J_k)^{1/d}
\gtrsim  
N(I \cap J_k)^{1/d}
=(N(I) N(J_k))^{1/d}
\gtrsim 
M_k^{2+\frac{\rho}{d}}, 
\end{align}
where the equality holds because $I$ and $J_k$ are relatively prime. 
Thus the elements of 
$\delta^{-1} I \cap (T^{-1} s + \delta^{-1} J)$ 
%are separated by the same amount.
are separated by 
$\gtrsim 
M_k^{2+\frac{\rho}{d}}.$ 
Since $T$ is a bounded linear isomorphism, 
the elements of $\mathcal{T}(I)$ 
are also separated by $\gtrsim M_k^{2+\frac{\rho}{d}}$. 
It follows that 
\begin{align*}
\sum_{t \in \mathcal{T}(I)} (1+\eta_k|t|)^{-50d}
&\lesssim
\sum_{t \in \ZZ^d} (1+\eta_k M_k^{2+\frac{\rho}{d}}|t|)^{-50d} 
\\
&\lesssim 
\left\{\begin{array}{cc}
    (M_k^{1+\tau} / M_k^{2+\frac{\rho}{d}})^d  & \text{if } \rho \leq d(\tau-1),  \\
    1 & \text{if } \rho > d(\tau-1). 
\end{array}\right.
\end{align*}
In any case, the sum is $\lesssim M_k^{d(\tau-1) -\rho}+1$. 
Plugging this estimate and \eqref{QMsize} into \eqref{h2hatsum2} 
completes the proof. 
\end{proof}

\end{document}
